% ====================================================================
% verification.tex (design note example)
% ====================================================================

\section{Verification Plan}
\label{sec:verification}

No boards have been built at this issue, so this section lists the tests that
will be carried out on the first ones, and the result each has to give for
the design to pass. The measured results will be added at the next issue.

\begin{table}[htbp]
\centering
\begin{tabularx}{\textwidth}{@{} l X X @{}}
\toprule
\textbf{Test} & \textbf{Method} & \textbf{Pass criterion} \\
\midrule
T1 & Apply DC levels of \qty{0.3}{\volt}, \qty{1.65}{\volt} and
     \qty{3.0}{\volt} to the input and measure the output.
   & Output within \qty{1}{\percent} of the input at each level. \\
\addlinespace
T2 & Apply a \qty{100}{\milli\volt} rms sine wave on a \qty{1.65}{\volt}
     offset, and step it from \qty{50}{\hertz} to \qty{500}{\hertz}.
   & Gain within \qty{1}{\decibel} of the DC gain at every step
     (requirement R4). \\
\addlinespace
T3 & Raise the same sine wave to \qty{1}{\volt} rms and measure the gain at
     \qty{9.5}{\kilo\hertz}, then at \qty{20}{\kilo\hertz},
     \qty{50}{\kilo\hertz} and \qty{100}{\kilo\hertz}.
   & At least \qty{35}{\decibel} of attenuation at \qty{9.5}{\kilo\hertz}
     (requirement R5), and no less than that at the higher frequencies. \\
\addlinespace
T4 & With the firmware running, toggle a spare pin in the interrupt routine
     and measure the period on an oscilloscope.
   & A period of \qty{100}{\micro\second}, with no visible jitter at
     \qty{1}{\micro\second} per division. \\
\bottomrule
\end{tabularx}
\caption{Tests planned for the first boards.}
\label{tab:tests}
\end{table}

Test T3 is taken beyond \qty{9.5}{\kilo\hertz} to find the point where the
attenuation stops increasing, as described at the end of
\cref{sec:performance}. If it turns out to fall below \qty{35}{\decibel}
within the range tested, the fix would be an op-amp with a wider bandwidth,
which is a change of part and not of the board.

\subsection{Open items}

\begin{itemize}
  \item The \qty{-40}{dBFS} limit on interference at the filter input, used
        in \cref{sec:requirements}, is an estimate. It should be measured on
        the first assembled unit, in its enclosure.
  \item The op-amp is specified by its parameters and not yet by a part
        number, which is to be fixed once stock and price have been checked.
\end{itemize}
